\documentclass[10pt,twocolumn,letterpaper]{article}

\usepackage[utf8]{inputenc}
\usepackage{graphicx}
\usepackage{amsmath}
\usepackage{amssymb}
\usepackage{booktabs}
\usepackage{hyperref}

\title{Project 5: High-Performance Semantic Correspondence via PEFT and Adaptive Windowing}
\author{Johnprice Osagie \and Mario Lapadula \and Giorgia Pugliese \and Riccardo Bellanca}
\date{April 2026}

\begin{document}

\maketitle

\begin{abstract}
This report details the implementation and evaluation of a high-performance pipeline for semantic correspondence using DINOv2 backbones. We investigate Parameter-Efficient Fine-Tuning (PEFT) strategies, specifically LoRA and BitFit, and introduce an Adaptive Window Soft-Argmax refinement stage. Our final model (LoRA + Adaptive Window) achieves a state-of-the-art result of 80.53\% PCK@0.1 on the SPair-71k test set, representing a significant improvement of +36\% over the zero-shot DINOv2 baseline.
\end{abstract}

\section{Introduction}
Semantic correspondence aims to establish dense matches between images of different objects within the same category. Challenges include viewpoint changes, scale variations, and inter-instance appearance shifts. We leverage the powerful features of DINOv2 and optimize the trade-off between performance and computational efficiency using PEFT.

\section{Methodology}
\subsection{Backbone and Extraction}
We utilize DINOv2 (ViT-B/14) as our feature extractor. Preliminary analysis confirms that semantic information is most concentrated in the deeper layers (Layer 11 for our pipeline).

\subsection{Efficient Fine-Tuning (LoRA vs. BitFit)}
We compare two PEFT strategies:
\begin{itemize}
    \item \textbf{LoRA}: Injects low-rank matrices into the attention layers. It provides the best accuracy (81.86\% on validation).
    \item \textbf{BitFit}: Optimizes only the bias parameters. It is extremely efficient and reaches competitive results (80.19\% on validation).
\end{itemize}

\subsection{Stage 3: Adaptive Window Refinement}
To improve sub-pixel precision, we implement an entropy-based \textbf{Adaptive Window Soft-Argmax}. This stage dynamically restricts the correspondence search area around the peak similarity, significantly boosting PCK@0.05 metrics.

\section{Experimental Results}
\subsection{Quantitative Analysis}
Table \ref{tab:pck} summarizes the results on the SPair-71k test set.

\begin{table}[h]
\centering
\caption{PCK results on SPair-71k Test Set.}
\label{tab:pck}
\begin{tabular}{@{}lcc@{}}
\toprule
Model & PCK@0.10 & PCK@0.05 \\ \midrule
Baseline (DINOv2) & 44.47\% & 24.65\% \\
BitFit (Smooth) & 75.40\% & 50.81\% \\
BitFit + Adaptive Window & 77.83\% & 55.48\% \\
LoRA (Smooth) & 78.13\% & 55.40\% \\
\textbf{LoRA + Adaptive Window} & \textbf{80.53\%} & \textbf{61.36\%} \\ \bottomrule
\end{tabular}
\end{table}

\subsection{Training Dynamics}
Figure \ref{fig:curves} shows the validation PCK@0.1 over 5 epochs. BitFit shows rapid convergence, while LoRA maintains a consistent accuracy lead.

\begin{figure}[h]
\centering
\includegraphics[width=\linewidth]{training_comparison.png}
\caption{Training curves for LoRA and BitFit.}
\label{fig:curves}
\end{figure}

\section{Discussion and Conclusion}
The integration of Adaptive Windowing proves critical, providing a +5.96\% boost to the PCK@0.05 score for LoRA. Our results demonstrate that even with minimal parameter updates (1.03\% trainable parameters), DINOv2 can be adapted to become a highly precise semantic matcher.

\end{document}
